% from ega4-20-fr.tex
\oldpage[IV-4]{252}
\subsection{Fonctions méromorphes relatives}
\EGAReaderTarget{ega4.c04.subsec.20-6}{IV.20.6}

\begin{env}[20.6.1]
\EGAReaderTarget{ega4.c04.statement.20-6-1}{IV.20.6.1}
Soient $S$ un préschéma, $X$ un $S$-préschéma qui est \emph{plat et localement de présentation finie} sur $S$ ; pour tout $s\in S$, nous noterons $X_s$ la fibre au point $s$ du morphisme structural $X\to S$.
En général, si $\varphi$ est une fonction méromorphe sur $X$, il n'est pas possible de lui associer de façon « naturelle », pour chaque $s\in S$, une fonction méromorphe « induite » sur $X_s$ \sref{IV.20.1.11}.
Pour tout ouvert $U$ de $X$, notons $T_{X/S}(U)$ l'ensemble des sections $t\in\Gamma(U,\sh{O}_X)$ telles que, pour tout $s\in S$, l'image $t_s$ de $t$ par l'homomorphisme canonique $\Gamma(U,\sh{O}_X)\to\Gamma(U\cap X_s,\sh{O}_{X_s})$ soit une section \emph{régulière} ; cela implique d'ailleurs, par l'équivalence de a) et b) dans \sref{IV.11.3.7}, que $t$ est elle-même une section \emph{régulière}.
Il est clair que $U\mapsto T_{X/S}(U)$ est un \emph{sous-faisceau} du faisceau d'ensemble $\sh{S}=\sh{S}(\sh{O}_X)$ (notation de \sref{IV.20.1.3}), que l'on note $\sh{T}_{X/S}$ ; on pose
\[
\label{IV.20.6.1.1}
  \sh{M}_{X/S}=\sh{O}_X[\sh{T}_{X/S}^{-1}]
  \tag{20.6.1.1}
\]
(notation de \sref{IV.20.1.1}) et on dit que ce faisceau d'anneaux est le \emph{faisceau des germes de fonctions méromorphes sur $X$ relatives à $S$} ; ses sections au-dessus de $X$ sont appelées les \emph{fonctions méromorphes sur $X$ relatives à $S$} et leur ensemble est noté $M(X/S)$.
Il est clair que $\sh{M}_{X/S}$ est un \emph{sous-faisceau} de $\sh{M}_X=\sh{O}_X[\sh{S}^{-1}]$ ; pour toute fonction méromorphe $\varphi\in M(X/S)$ et tout $s\in S$, l'image réciproque de $\varphi$ par le morphisme injection canonique $j_s:X_s\to X$ est alors définie \sref{IV.20.1.11}, et notée $\varphi_s$.
\end{env}

\begin{env}[20.6.2]
\EGAReaderTarget{ega4.c04.statement.20-6-2}{IV.20.6.2}
Soit maintenant $\sh{F}$ un $\sh{O}_X$-Module quasi-cohérent ; on pose
\[
\label{IV.20.6.2.1}
  \sh{M}_{X/S}(\sh{F})=\sh{F}[\sh{T}_{X/S}^{-1}]=\sh{F}\otimes_{\sh{O}_X}\sh{M}_{X/S};
  \tag{20.6.2.1}
\]
les sections de $\sh{M}_{X/S}(\sh{F})$ sont dites \emph{sections méromorphes de $\sh{F}$ au-dessus de $X$, relatives à $S$} et leur ensemble est noté $M(X/S,F)$.
L'homomorphisme canonique $\sh{F}\to\sh{M}_{X/S}(\sh{F})$ n'est pas nécessairement injectif ; lorsqu'il l'est, on dit que $\sh{F}$ est \emph{sans torsion relativement à $S$} : cela signifie que pour tout ouvert $U$ de $X$ et toute section $t\in T_{X/S}(U)$, $t$ est $(\sh{F}|U)$-régulière ; cette condition est remplie \emph{a fortiori} lorsque $\sh{F}$ est \emph{strictement sans torsion} \sref{IV.20.1.5}.
Dans ce dernier cas, il résulte aussitôt des définitions \sref{IV.20.1.2} que l'homomorphisme canonique de $\sh{O}_X$-Modules
\[
\label{IV.20.6.2.2}
  \sh{M}_{X/S}(\sh{F})\to\sh{M}_X(\sh{F})
  \tag{20.6.2.2}
\]
est \emph{injectif}, de sorte que les sections méromorphes de $\sh{F}$ relatives à $S$ sont des sections méromorphes de $\sh{F}$ au sens de \sref{IV.20.1.3}.
\end{env}

\begin{proposition}[20.6.3]
\label{IV.20.6.3}
L'image par l'homomorphisme injectif \sref{IV.20.2.10.1} du sous-faisceau $\sh{M}_{X/S}$ de $\sh{M}_X$ est le sous-faisceau $\sh{M}'_{X/S}$ des pseudo-fonctions sur $X$ relatives à $S$ (i.e.\ des pseudo-fonctions dont le domaine de définition est universellement schématiquement dense relativement à $S$ \sref{IV.20.5.4}).
\end{proposition}

\begin{proof}
On peut évidemment se borner à prouver que l'image de $M(X/S)$ par l'homomorphisme canonique $M(X)\to M'(X)$ est égale à $M'(X/S)$ ; la proposition est alors conséquence de la proposition plus générale suivante :
\end{proof}

